\section{Critical Reflection}
% ─────────────────────────────────────────────────────────────

This project has provided a challenging and illuminating end-to-end experience with applied machine learning—from raw multi-gigabyte network captures to production-oriented model comparison. The following reflection distils the most significant technical and epistemic lessons acquired.

\textbf{The Gap Between Theory and Practice.}
The formal guarantees of classical anomaly detectors—LOF's density consistency, OC-SVM's maximum-margin optimality—are formulated in idealised settings that frequently do not hold in practice. Our results starkly illustrate this: LOF, despite a well-studied theoretical foundation, achieves near-random AUC (0.44) on a 70-feature dataset. The culprit is not the algorithm's design but the fundamental statistical phenomenon of distance concentration in high-dimensional spaces, a phenomenon that is theoretically understood \cite{zimek2012survey} yet routinely underestimated in introductory ML curricula. This experience reinforced that selecting a model must be driven by the geometry of the data, not by pedagogical familiarity.

\textbf{The Importance of Hyperparameter Sensitivity Analysis.}
The hyperparameter tuning exercise revealed that the contamination parameter is far more impactful than the number of trees or neighbours. Small misspecifications of contamination—say, setting it to 0.20 when the true rate is 0.15—can cost 2–4 percentage points of F1 on imbalanced data. This sensitivity is not prominently communicated in the original algorithm papers, which typically evaluate on balanced benchmarks. Systematic grid search, even at modest computational cost, is therefore an essential step rather than an optional refinement.

\textbf{On Evaluation Design.}
Perhaps the most important methodological lesson concerns evaluation. Standard accuracy on an 80/20 imbalanced dataset is almost entirely uninformative. A classifier predicting ``normal'' for every input achieves 80\% accuracy—yet catches zero attacks. Our three-perspective framework (original distribution, balanced subset, per-class analysis) was indispensable for exposing real model behaviour. The ensemble autoencoder's 96.3\% recall at 57\% precision illustrates the canonical cybersecurity trade-off: the operational cost of false alarms is tractable; the cost of a missed attack is potentially catastrophic. Designing evaluation protocols that reflect deployment objectives—rather than publishing convenience—is a core competency that this project helped develop.

\textbf{Ethical Dimensions.}
Anomaly detection systems do not operate in a vacuum. False positives can flag legitimate users as threats, with potential legal and professional consequences. False negatives can enable real harm. Deploying such systems in sensitive infrastructure—healthcare networks, critical national infrastructure—demands transparency about error rates and careful threshold calibration informed by domain expertise. Our models were trained on a dataset collected from a Canadian university network in 2017; their generalisation to other network topologies, cultural contexts, and evolving attack vectors cannot be assumed. Responsible deployment would require continuous monitoring, regular retraining, and explicit human oversight of high-risk predictions.

\textbf{Broader Takeaways.}
This project bridged the gap between the mathematical abstractions of the classroom and the messy realities of engineering ML systems: handling out-of-memory errors through stratified sampling, preventing data leakage by fitting scalers exclusively on training splits, and balancing the computational feasibility of $\mathcal{O}(n^2)$ algorithms against their theoretical appeal. These engineering decisions are invisible in published papers yet dominate practical project time. Learning to make them thoughtfully—and to document the rationale—is, ultimately, as valuable as learning the algorithms themselves.
